% ==================================================
% Section 2
% ==================================================
\section{Model Reference}

The objective of the adaptive controller is to make the plant output \(y(t)\) track the output of a stable reference model \(y_m(t)\). For the third-order system, the reference model is chosen as:
\begin{equation}
y_m^{(3)}
+\alpha_2\ddot{y}_m
+\alpha_1\dot{y}_m
+\alpha_0 y_m
=
\alpha_0 r(t)
\label{eq:reference_model}
\end{equation}

where \(r(t)\) is the reference input and the constants \(\alpha_0,\alpha_1,\alpha_2>0\) are chosen such that the reference model is stable.

\subsection{System Parametrization}

The plant model is:
\begin{equation}
y^{(3)}
+a_1y''
+a_2y'
+a_3y
+b_1(y'')^3
+b_2(y')^3
+b_3\sin(y)
=
cu.
\end{equation}
Rearranging gives:
\begin{equation}
y^{(3)}
=
-\theta^T\phi
+
cu,
\label{eq:plant_parametric}
\end{equation}

where the unknown parameter vector is:
\begin{equation}
\theta =
\begin{bmatrix}
a_1 & a_2 & a_3 & b_1 & b_2 & b_3
\end{bmatrix}^{T},
\end{equation}

and the known regression vector is:
\begin{equation}
\phi =
\begin{bmatrix}
y'' & y' & y & (y'')^3 & (y')^3 & \sin(y)
\end{bmatrix}^{T}.
\end{equation}
The parameters \(a_i\) and \(b_i\) are assumed to be unknown constants, while the input gain \(c\) is assumed to be known and non-zero.

\subsection{Tracking Error and Filtered Error}

The tracking error is defined as:
\begin{equation}
e = y-y_m.
\end{equation}
Following the adaptive-control design for nonlinear systems of any order, a filtered tracking error is defined. Since the plant is third-order, the filtered error is chosen as:
\begin{equation}
z
=
\ddot{e}
+
\lambda_2\dot{e}
+
\lambda_1 e
\label{eq:filtered_error}
\end{equation}

where \(\lambda_1,\lambda_2>0\) are design parameters, such as that the polynomial
$
s^2+\lambda_2s+\lambda_1
$
is selected to be stable.
\\

\noindent The filtered error \(z\) combines the tracking error, velocity error, and acceleration error into one scalar variable. If \(z(t)\rightarrow0\), then the tracking error follows a stable second-order error equation and \(e(t)\rightarrow0\).

\subsection{Filtered Error Dynamics}

From the definition of \(z\), its derivative is:
\begin{equation}
\dot{z}
=
e^{(3)}
+
\lambda_2\ddot{e}
+
\lambda_1\dot{e}
\end{equation}

since
\begin{equation}
e^{(3)}
=
y^{(3)}
-
y_m^{(3)}
\end{equation}

we obtain:
\begin{equation}
\dot{z}
=
y^{(3)}
-
y_m^{(3)}
+
\lambda_2\ddot{e}
+
\lambda_1\dot{e}.
\end{equation}
To simplify the controller design, an auxiliary reference signal is defined as:
\begin{equation}
y_r^{(3)}
=
y_m^{(3)}
-
\lambda_2\ddot{e}
-
\lambda_1\dot{e}
\label{eq:yr_definition}
\end{equation}
Substituting \eqref{eq:yr_definition} into the expression for \(\dot{z}\) gives:
\begin{equation}
\dot{z}
=
y^{(3)}
-
y_r^{(3)}.
\label{eq:z_dot_compact}
\end{equation}
This representation expresses the filtered error dynamics in terms of the plant output derivative \(y^{(3)}\) and the auxiliary reference signal \(y_r^{(3)}\). The resulting form will be used in the next section to construct the control law and derive the desired closed-loop error dynamics.

\subsection{Adaptive Control Law}

In order to force the filtered error \(z(t)\) to converge to zero, the desired filtered error dynamics are selected as:
\begin{equation}
\dot{z}
=
-k_z z
\label{eq:desired_z_dynamics}
\end{equation}

where \(k_z>0\) is a design parameter.
\\

\noindent Equation \eqref{eq:desired_z_dynamics} is a stable first-order differential equation whose solution is:
\begin{equation}
z(t)
=
z(0)e^{-k_z t}
\end{equation}
Therefore, the filtered error converges exponentially to zero. The control law will now be designed such that the closed-loop system follows the desired dynamics \eqref{eq:desired_z_dynamics}.
\noindent First, we assume that the parameter vector \(\theta\) is known. In this case, the control input is chosen as:
\begin{equation}
u
=
\frac{1}{c}
\left[
\theta^T\phi
+
y_r^{(3)}
-
k_z z
\right]
\label{eq:nominal_control_law}
\end{equation}
Substituting \eqref{eq:nominal_control_law} into the plant model \eqref{eq:plant_parametric} gives:
\begin{align}
y^{(3)}
&=
-\theta^T\phi
+
c
\left(
\frac{1}{c}
\left[
\theta^T\phi
+
y_r^{(3)}
-
k_z z
\right]
\right) \\
&=
-\cancel{\theta^T\phi}
+
\cancel{\theta^T\phi}
+
y_r^{(3)}
-
k_z z \\
&=
y_r^{(3)}
-
k_z z
\end{align}
In this case, we can see that the filtered error dynamics become identical to \eqref{eq:z_dot_compact} as we assumed:
\begin{equation}
y^{(3)} - y_r^{(3)}
=
\dot{z}
=
-k_z z.
\end{equation}
Therefore, the closed-loop system achieves the desired filtered error dynamics and the filtered error converges exponentially to zero.

\subsubsection*{Adaptive Controller}

In practice, the parameters \(a_i\) and \(b_i\) are unknown. Therefore, the true parameter vector \(\theta\) cannot be used directly and must be replaced by its estimate \(\hat{\theta}\). Using the definition of \(y_r^{(3)}\) from \eqref{eq:yr_definition}, the controller can also be written as:
\begin{equation}
u
=
\frac{1}{c}
\left[
\hat{\theta}^{T}\phi
+
y_m^{(3)}
-
\lambda_2\ddot{e}
-
\lambda_1\dot{e}
-
k_z z
\right].
\end{equation}
We define the parameter estimation error as:
\begin{equation}
\tilde{\theta}
=
\hat{\theta}
-
\theta.
\end{equation}
Substituting \eqref{eq:nominal_control_law} into the plant model \eqref{eq:plant_parametric} yields:
\begin{align}
y^{(3)}
&=
-\theta^T\phi
+
\hat{\theta}^{T}\phi
+
y_r^{(3)}
-
k_z z \\
&=
y_r^{(3)}
-
k_z z
+
\tilde{\theta}^{T}\phi.
\label{eq:adaptive_control_law}
\end{align}
Using \eqref{eq:z_dot_compact}, the filtered error dynamics become:
\begin{equation}
y^{(3)} - y_r^{(3)}
=
\dot{z}
=
-k_z z
+
\tilde{\theta}^{T}\phi.
\label{eq:z_dynamics}
\end{equation}
Compared to the nominal case, an additional term \(\tilde{\theta}^{T}\phi\) appears due to parameter uncertainty. The purpose of the adaptation law is to update \(\hat{\theta}\) online and compensate for this uncertainty while preserving the desired stability properties.

\subsection{Adaptation Law}

The parameter estimates are updated online using the adaptation law:
\begin{equation}
\dot{\hat{\theta}}
=
-\Gamma \phi z,
\label{eq:adaptation_law}
\end{equation}

where \(\Gamma>0\) is a positive definite adaptation gain matrix.
\\

\noindent The matrix \(\Gamma\) determines the adaptation speed. Larger values of \(\Gamma\) generally lead to faster parameter adaptation, but may also increase sensitivity to measurement noise and modeling uncertainties.

\subsection{Lyapunov Stability Proof}

To prove stability, we consider the Lyapunov candidate:
\begin{equation}
V
=
\frac{1}{2}z^2
+
\frac{1}{2}\tilde{\theta}^{T}\Gamma^{-1}\tilde{\theta}.
\label{eq:lyapunov_candidate}
\end{equation}
The first term represents the filtered tracking error energy, and the second term represents the parameter estimation error energy. Since \(\Gamma>0\), the function \(V\) is positive definite with respect to \(z\) and \(\tilde{\theta}\). Differentiating \eqref{eq:lyapunov_candidate} gives:
\begin{equation}
\dot{V}
=
z\dot{z}
+
\tilde{\theta}^{T}\Gamma^{-1}\dot{\tilde{\theta}}.
\end{equation}
Since the real parameters are constant:
\begin{equation}
\dot{\tilde{\theta}}
=
\dot{\hat{\theta}}.
\end{equation}
Substituting \eqref{eq:z_dynamics} and \eqref{eq:adaptation_law} gives:
\begin{align}
\dot{V}
&=
z
\left(
-k_z z
+
\tilde{\theta}^{T}\phi
\right)
+
\tilde{\theta}^{T}\Gamma^{-1}
\left(
-\Gamma\phi z
\right)\\
&=
-k_z z^2
+
z\tilde{\theta}^{T}\phi
-
\tilde{\theta}^{T}\phi z.
\end{align}
The last two terms are equal scalars with opposite signs, and therefore they cancel:
\begin{equation}
\dot{V}
=
-k_z z^2
+
\cancel{z\tilde{\theta}^{T}\phi}
-
\cancel{\tilde{\theta}^{T}\phi z}
=
-k_z z^2.
\label{eq:lyapunov_derivative}
\end{equation}
 Since \(k_z>0\), we obtain:
\begin{equation}
\dot{V}
=
-k_z z^2
\leq 0.
\end{equation}
Therefore, the Lyapunov function never increases with time. This implies that \(z(t)\) and \(\tilde{\theta}(t)\) remain bounded. Furthermore, the filtered error \(z(t)\) converges to zero:
\begin{equation}
z(t)\rightarrow0
\qquad
\text{as}
\qquad
t\rightarrow\infty.
\end{equation}
 From the definition of \(z\):
\begin{equation}
z
=
\ddot{e}
+
\lambda_2\dot{e}
+
\lambda_1e.
\end{equation}
When \(z(t)\rightarrow0\), the tracking error satisfies:
\begin{equation}
\ddot{e}
+
\lambda_2\dot{e}
+
\lambda_1e
=
0.
\end{equation}
 Since \(\lambda_1,\lambda_2>0\) are chosen to give stable error dynamics, it follows that:
\begin{equation}
e(t)\rightarrow0,
\qquad
\dot{e}(t)\rightarrow0.
\end{equation}
Finally, since \(e=y-y_m\), convergence of the tracking error implies:
\begin{equation}
y(t)\rightarrow y_m(t).
\end{equation}
Thus, the proposed adaptive controller guarantees bounded closed-loop signals and asymptotic tracking of the reference model output. $\blacksquare$

\bigskip